\begin{figure}[htbp]
\centering
\begin{tikzpicture}[
  vertex/.style={circle,draw=black!65,fill=white,inner sep=2pt},
  edge/.style={-{Stealth[length=1.8mm]},draw=black!65},
  lab/.style={font=\small,fill=white,inner sep=1pt}]
 \node (in) at (0,0) {$M$};
 \node[vertex] (a) at (1.6,0) {$a_1$};
 \node[vertex] (b) at (4.7,0) {$a_2$};
 \node[vertex] (c) at (7.8,0) {$a_3$};
 \node (out) at (9.6,0) {$dq$};
 \draw[edge,line width=1pt] (in) -- (a);
 \draw[edge,line width=1pt] (a) -- node[lab,above] {$Q_1$} (b);
 \draw[edge,line width=1pt] (b) -- node[lab,above] {$Q_2$} (c);
 \draw[edge,line width=1pt] (c) -- node[lab,above] {suffix} (out);
 \node[draw=black!60,rounded corners=2pt,fill=black!3,
       minimum width=8.2cm,minimum height=8mm]
       (tensor) at (4.7,-1.4)
       {$U_i(a_1,a_2,a_3)$: pulled-back label tensor,
        $\|U_i\|_{\ell^2}\leq C$};
 \draw[edge] (a) -- (a |- tensor.north);
 \draw[edge] (b) -- (b |- tensor.north);
 \draw[edge] (c) -- (c |- tensor.north);
 \node[font=\small] at (4.7,-2.05)
      {nonempty label branches at every attachment};
\end{tikzpicture}
\caption{The derivative continuing into the later suffix defines
a marked path.  Side branches are absorbed into a label tensor
with bounded row square norm.  After branch-free segments have been
combined, every intervening $Q_r$ separates two nonempty sets of
label leaves.  Its off-diagonal contribution is small by
Cauchy--Schwarz and the row and column sums of the squared-entry
matrices.  Pinching all such segments ties the attachment indices
and returns the diagonal estimate for $M\circ E$.}
\label{fig:marked-continuation}
\end{figure}
